\documentclass{ximera}
\author{Jeffrey Kuan}
\title{Math 1050 Notes: Section 1.4}
\license{CC: 0}

\begin{document}

\begin{abstract}
    Section 1.4 --- Exponents and the Order of Operations Agreement.
    Worked examples, then practice problems with answers.
\end{abstract}

\maketitle

\textbf{Order of Operations Agreement.} Work in this order: \textbf{P}arentheses,
\textbf{E}xponents, \textbf{M}ultiplication and \textbf{D}ivision (left to right), then
\textbf{A}ddition and \textbf{S}ubtraction (left to right).

Watch the difference between $-2^{2}$ and $(-2)^{2}$: the exponent applies only to the 2
in the first, so $-2^{2}=-4$, while $(-2)^{2}=4$.

\section*{Evaluate each expression}

\begin{problem}
\textbf{1.} $3^{3}$ \qquad \textbf{Answer:} $\answer{27}$

\begin{hint}
\textbf{Solution.} $3^{3}=3\cdot 3\cdot 3=27$.
\end{hint}
\end{problem}

\begin{problem}
\textbf{2.} $(-3)^{4}$ \qquad \textbf{Answer:} $\answer{81}$

\begin{hint}
\textbf{Solution.} The exponent is even, so the result is positive: $(-3)^{4}=81$.
\end{hint}
\end{problem}

\begin{problem}
\textbf{3.} $-2^{2}$ \qquad \textbf{Answer:} $\answer{-4}$

\begin{hint}
\textbf{Solution.} The exponent applies only to the 2, so $-2^{2}=-(2\cdot 2)=-4$.
\end{hint}
\end{problem}

\begin{problem}
\textbf{4.} $\left(\dfrac{2}{3}\right)^{2}$ \qquad \textbf{Answer:} $\answer{\frac{4}{9}}$

\begin{hint}
\textbf{Solution.} $\dfrac{2}{3}\cdot\dfrac{2}{3}=\dfrac{2^{2}}{3^{2}}=\dfrac{4}{9}$.
\end{hint}
\end{problem}

\begin{problem}
\textbf{5.} $\dfrac{2^{2}}{3}$ \qquad \textbf{Answer:} $\answer{\frac{4}{3}}$
\end{problem}

\begin{problem}
\textbf{6.} $\left(-\dfrac{3}{4}\right)^{2}$ \qquad
\textbf{Answer:} $\answer{\frac{9}{16}}$

\begin{hint}
\textbf{Solution.} $\dfrac{(-3)^{2}}{4^{2}}=\dfrac{9}{16}$.
\end{hint}
\end{problem}

\section*{Simplify using the Order of Operations Agreement}

\begin{problem}
\textbf{8.} $14-2^{2}-|4-7|$ \qquad \textbf{Answer:} $\answer{7}$

\begin{hint}
\textbf{Solution.} $14-4-3=10-3=7$.
\end{hint}
\end{problem}

\begin{problem}
\textbf{9.} $24 \div \dfrac{3^{2}}{8-5}-(-5)$ \qquad \textbf{Answer:} $\answer{13}$

\begin{hint}
\textbf{Solution.} $24 \div \dfrac{9}{3}+5 = 24 \div 3+5 = 8+5 = 13$.
\end{hint}
\end{problem}

\begin{problem}
\textbf{10.} $\left(\dfrac{3}{4}\right)^{2}-\left(\dfrac{1}{2}\right)^{3} \div \dfrac{3}{5}$
\qquad \textbf{Answer:} $\answer{\frac{17}{48}}$

\begin{hint}
\textbf{Solution.} $\dfrac{1}{8}\div\dfrac{3}{5}=\dfrac{1}{8}\cdot\dfrac{5}{3}=\dfrac{5}{24}$,
so the expression is
$\dfrac{9}{16}-\dfrac{5}{24}=\dfrac{27}{48}-\dfrac{10}{48}=\dfrac{17}{48}$.
\end{hint}
\end{problem}

\begin{problem}
\textbf{11.} $3\cdot(15-2\cdot 3)-36 \div 3$ \qquad \textbf{Answer:} $\answer{15}$

\begin{hint}
\textbf{Solution.} $3(15-6)-12 = 3\cdot 9-12 = 27-12 = 15$.
\end{hint}
\end{problem}

\section*{Your turn}

\begin{problem}
\textbf{1.} $\left(-\dfrac{2}{5}\right)^{2}$ \qquad \textbf{Answer:} $\answer{\frac{4}{25}}$
\end{problem}

\begin{problem}
\textbf{2.} $(-5)^{2}$ \qquad \textbf{Answer:} $\answer{25}$
\end{problem}

\begin{problem}
\textbf{3.} $\left(\dfrac{1}{2}\right)^{3} \cdot 8$ \qquad \textbf{Answer:} $\answer{1}$
\end{problem}

\begin{problem}
\textbf{4.} $(0.5)^{2} \cdot 3^{3}$ \qquad \textbf{Answer:} $\answer{6.75}$
\end{problem}

\begin{problem}
\textbf{5.} $3^{2} \cdot 2^{2}-3$ \qquad \textbf{Answer:} $\answer{33}$

\begin{hint}
\textbf{Solution.} $9\cdot 4-3=36-3=33$.
\end{hint}
\end{problem}

\begin{problem}
\textbf{6.} $9-4+2 \cdot 2$ \qquad \textbf{Answer:} $\answer{9}$

\begin{hint}
\textbf{Solution.} $9-4+4=5+4=9$.
\end{hint}
\end{problem}

\begin{problem}
\textbf{7.} $-2^{2}+4\bigl(16 \div (3-5)\bigr)$ \qquad \textbf{Answer:} $\answer{-36}$

\begin{hint}
\textbf{Solution.} $-4+4(16 \div (-2)) = -4+4(-8) = -4-32 = -36$.
\end{hint}
\end{problem}

\begin{problem}
\textbf{8.} $16-4 \cdot \dfrac{3^{3}-7}{2^{3}+2}-(-2)^{2}$ \qquad
\textbf{Answer:} $\answer{4}$

\begin{hint}
\textbf{Solution.} $16-4\cdot\dfrac{20}{10}-4 = 16-8-4 = 4$.
\end{hint}
\end{problem}

\end{document}
